\clearpage
\chapter*{РЕФЕРАТ}
\addcontentsline{toc}{chapter}{РЕФЕРАТ}
Выпускная квалификационная работа посвящена разработке вычислительной модели и программного комплекса для расчёта термодинамических свойств воды и водяного пара.
В качестве нормативной основы используется промышленный стандарт IAPWS-IF97 в редакции 2012 года \cite{iapws_if97_2012}.

Цель работы: разработать единое вычислительное ядро и набор каналов доставки результатов, обеспечивающих как интерактивные инженерные расчёты, так и интеграцию в сервисную инфраструктуру.

Задачи работы включают реализацию вычислительного ядра на языке Rust по стандарту IAPWS-IF97; поддержку прямых и обратных маршрутов расчёта (в том числе по парам $pT$, $pH$, $pS$, $pX$); численное расширение маршрута расчёта по паре $(\rho, T)$ на всей рабочей области на основе метода секущих; разработку настольного приложения на FLTK и desktop-оболочки на базе Yew, WebAssembly и Tauri; реализацию высокопроизводительного gRPC-сервиса для пакетных расчётов; контейнеризацию и подготовку инфраструктуры развёртывания в Kubernetes; автоматизацию сборки, тестирования и публикации артефактов через CI/CD.

Объект исследования: термодинамические свойства воды и водяного пара.
Предмет исследования: алгоритмы и программная реализация вычисления свойств в области применимости IAPWS-IF97, а также архитектурные решения, обеспечивающие повторное использование вычислительного ядра в различных приложениях.

Методы исследования: аналитические зависимости IAPWS-IF97, численные методы решения нелинейных уравнений (метод секущих), модульное проектирование, валидация диапазонов входных параметров и тестирование по контрольным данным.

Результаты работы: реализовано вычислительное ядро \texttt{if97\_core} с единым API, типобезопасными единицами измерения и унифицированной структурой результата; разработаны два пользовательских приложения (нативное и web-технологическое); реализован gRPC-сервис с разделением I/O и CPU-нагрузки и механизмами ограничения нагрузки; подготовлены Docker-образы, Kubernetes-манифесты и сценарии CI/CD.

Практическая значимость работы заключается в возможности использовать единый набор вычислений в инженерном рабочем месте (desktop), в составе кроссплатформенного интерфейса (Yew/Tauri) и как сервис для пакетных расчётов и интеграции в другие системы.

Работа состоит из введения, пяти глав, заключения, списка использованных источников и приложений.

\textbf{Ключевые слова:}\\
вода, водяной пар, термодинамические свойства, IAPWS-IF97, численные методы, метод секущих;\\
Rust, FLTK, WebAssembly, Tauri;\\
gRPC, Docker, Kubernetes, CI/CD.
